\scanpage{571}
\mspage{589}
% source-position: top lines 1--4
is then unique, and compatible with \(i_{\mathcal{G},G}\)
\inserted{\(i_{Z,G}\)}\unsure{;}
\deleted{\unsure{transforms}}\inserted{maps}
\(\deleted{\unsure{N}}\inserted{Z}_\rho^{\mathcal{G}}\) into
\(\deleted{\unsure{N}}\inserted{Z}_\rho^{\mathcal{G}'}\),
\(\deleted{\unsure{N}}\inserted{Z}_\sigma^{\mathcal{G}}\) into
\(\deleted{\unsure{N}}\inserted{Z}_\sigma^{\mathcal{G}'}\),
\unsure{which} extends to
\(\varphi_\rho:\diagram{\begin{tikzcd}N_\rho^{\mathcal{G}}\arrow[r]&N_\rho^{\mathcal{G}'}\end{tikzcd}}\)
inducing
\(\diagram{\begin{tikzcd}N_\rho^{*\mathcal{G}}\arrow[r]&N_\rho^{*\mathcal{G}'}\end{tikzcd}}\),
and extends to
\(\varphi_\sigma:\diagram{\begin{tikzcd}N_\sigma^{\mathcal{G}}\arrow[r]&N_\sigma^{\mathcal{G}'}\end{tikzcd}}\)
inducing
\(\diagram{\begin{tikzcd}N_\sigma^{*\mathcal{G}}\arrow[r]&N_\sigma^{*\mathcal{G}'}\end{tikzcd}}\)\ldots)

% source-position: upper lines 5--7
\deleted{\unsure{it} \illegible{4}}
\deleted{d) \(\exists\) homomorphism of groups}

We shall then say that \(\varphi\) is ``extendable''.

% source-position: upper-middle lines 8--10
\uline{Corollary}\quad We can then extend \(\varphi\) to
\[
(93)\qquad\varphi_M:
\diagram{\begin{tikzcd}
S_0G_{\rho,\sigma}^{\mathcal{G}}\arrow[r]&S_0G_{\rho,\sigma}^{\mathcal{G}'}
\end{tikzcd}},
\]
compatible with \(\Theta_Z,\Theta_\rho,\Theta_\sigma,\Theta_G\)
(taking into account \(\varphi_Z,\varphi_\rho,\varphi_\sigma,\varphi_G\)),
and this in a way ``compatible with \unsure{everything}'', in particular
with diagrams (45) (62).

% source-position: middle-lower lines 11--14
\uline{Example}\quad Suppose
\(\Sigma_{\rho,\sigma}^{\mathcal{G}'}=\Sigma_{\rho,\sigma}^{\mathcal{G}'\,\mathrm{eff}}\),
then the conditions of the proposition are \unsure{automatically}
satisfied (whatever the surjective homomorphism
\(\diagram{\begin{tikzcd}\mathcal{G}\arrow[r]&\mathcal{G}'\end{tikzcd}}\)).

% source-position: lower lines 15--19
\uline{Remark}\quad It may be that we are given a subgroup
\(Z_{\rho,\sigma}^{\unsure{\natural|\sharp}}\subset Z_{\rho,\sigma}^{\mathcal{G}}\),
and that we consider \deleted{\unsure{a)}} condition a) of the proposition
only for \(u\in Z_{\rho,\sigma}^{\unsure{\natural|\sharp}}\).
But with the help of \(Z_{\rho,\sigma}^{\unsure{\natural|\sharp}}\), we construct
\margin{We must suppose \unsure{in addition} that
\(\tau'_Z\in Z_{\rho,\sigma}^{*}\)}
\[
(94)\quad
G_{\rho,\sigma}^{\unsure{\natural|\sharp}}=
Z_{\rho,\sigma}^{\unsure{\natural|\sharp}}\cdot\mathcal{G}^+,
\quad Z_\rho^{\unsure{\natural|\sharp}}=Z_\rho\cap G_{\rho,\sigma}^{\unsure{\natural|\sharp}},
\quad Z_\sigma^{\unsure{\natural|\sharp}}=Z_\sigma\cap G_{\rho,\sigma}^{\unsure{\natural|\sharp}}
\]
% source-position: bottom lines 20--22
whence \(N_\rho^{\unsure{\natural|\sharp}}\deleted{_{\unsure{\rho,\sigma}}}\)
and \(N_\rho^{\unsure{\natural|\sharp}*}\deleted{_{\unsure{\rho,\sigma}}}\),
\(N_\sigma^{\unsure{\natural|\sharp}}\) and
\(N_\sigma^{\unsure{\natural|\sharp}*}\),
\(\gamma_{\rho,\sigma}^{\unsure{\natural|\sharp}}\), then
\(\unsure{S}_0M_{\rho,\sigma}^{\unsure{\natural|\sharp}}\) and
\(M_{\rho,\sigma}^{\unsure{\natural|\sharp}}\).
\deleted{We have a variant of \unsure{the}}
This being so, we have an obvious variant of the
\mspage{570}

% Printed archive footer outside the manuscript frame (not manuscript):
% 570
% https://grothendieck.umontpellier.fr
% Université de Montpellier, UMR CNRS 5149, Case courrier 051, Place Eugène Bataillon
% 34095 Montpellier cedex 5 - France

\scanpage{572}
% source-position: top lines 1--2
proposition, giving conditions for \unsure{extension|extension} of \(\varphi\) to
\(\diagram{\begin{tikzcd}
S_0M_{\rho,\sigma}^{*}\arrow[r]&S_0M_{\rho,\sigma}^{\mathcal{G}'}
\end{tikzcd}}\).

% source-position: upper abandoned block, lines 3--7
\deleted{\uline{Universal case}
\[
\mathcal{G}=\operatorname{Gl}(2,\mathbb Z)^\wedge,\quad
\rho=\begin{pmatrix}0&1\\-1&1\end{pmatrix},\quad
\sigma=\begin{pmatrix}0&-1\\1&0\end{pmatrix}
\]
\[
\tau'=\begin{pmatrix}0&1\\1&0\end{pmatrix},\quad
\unsure{\{}\ \tau=\begin{pmatrix}-1&0\\0&+1\end{pmatrix};\quad
\varepsilon_0=\begin{pmatrix}1&-1\\0&1\end{pmatrix},
\]
\deleted{We have}
\[
\varepsilon_1=\begin{pmatrix}\unsure{1}&0\\1&1\end{pmatrix}
\deleted{^{\unsure{\varepsilon}}}.
\]
Since \(\rho\)}
\mspage{571}

% Printed archive footer outside the manuscript frame (not manuscript):
% 571
% https://grothendieck.umontpellier.fr
% Université de Montpellier, UMR CNRS 5149, Case courrier 051, Place Eugène Bataillon
% 34095 Montpellier cedex 5 - France

\scanpage{573}
\mspage{590}
% source: upper portion, heading, lines 1--5, (1)--(2).

$8^\circ)$ \uline{The universal case} $\mathcal{G}=\mathrm{Gl}(2,\mathbb{Z})^{\wedge}$.

\unsure{.} Let us consider the exact sequence

\diagram{\begin{tikzcd}
\text{(1)}\quad 1 \arrow[r] & \{\pm1\} \arrow[r] & \mathcal{N}_{1,1}^{\sim} \arrow[r] \arrow[d,equal] & \mathcal{M}_{0,3}^{\sim} \arrow[r] & 1 \\
 & & \mathcal{M} & &
\end{tikzcd}}

\deleted{the group $\{\pm1\}$}
We have
\[
(2)\qquad \{\pm1\}=\unsure{\mathrm{Centr}|\mathrm{Centre}}(\mathcal{M})=\unsure{\mathrm{Centr}|\mathrm{Centre}}(\widehat{\mathcal{G}}^{+}),
\]

% source: oblique left margin beside (2)--(3), four identifications.
\margin{\unsure{we have taken|we have set}
\[
\begin{aligned}
\mathcal{G}&=\unsure{\mathcal{T}_{1,1}|\mathcal{G}_{1,1}}\simeq\mathrm{Gl}(2,\mathbb{Z}),\\
\widehat{\mathcal{G}}&=\widehat{\mathcal{T}}_{1,1}\simeq\mathrm{Gl}(2,\mathbb{Z})^{\wedge},\\
\unsure{\widehat{\mathcal{G}}^{+}|\mathcal{G}^{+}}&=\mathcal{T}_{1,1}^{+}=\mathrm{Sl}(2,\mathbb{Z}),\\
\widehat{\mathcal{G}}^{+}&=\mathrm{Sl}(2,\mathbb{Z})^{\wedge}.
\end{aligned}
\]}

% source: middle portion, lines 6--12, (3)--(4).
whence by (1) a map

\diagram{\begin{tikzcd}
\text{(3)}\quad \mathcal{M}_{0,3}^{\sim} \arrow[r] & \mathrm{Aut}(\mathcal{M}) \arrow[r,"{\inserted{\substack{\text{restriction to}\\ \mathcal{G}^{+\wedge}=\mathrm{Sl}(2,\mathbb{Z})^{\wedge}}}}"] & \mathrm{Aut}(\widehat{\mathcal{G}}^{+\deleted{\unsure{\wedge}}})
\end{tikzcd}}

\deleted{in the case}
Besides, \unsure{every} \unsure{automorphism} of $\widehat{\mathcal{G}}^{+}$ \unsure{preserves}
its center, hence \unsure{passes} to the quotient, whence

\diagram{\begin{tikzcd}
\text{(4)}\quad \mathrm{Aut}(\widehat{\mathcal{G}}^{+}) \arrow[r] & \mathrm{Aut}(\widehat{\mathcal{G}}^{+}/\!\pm1\simeq\deleted{\unsure{\mathcal{M}}}\mathcal{G}_{0,3}^{+\wedge})
\end{tikzcd}}

% source: lower-middle portion, lines 13--16, triangle (5).
and \unsure{the composite|the composite} of (3) with (4) is the canonical
homomorphism:

\diagram{\begin{tikzcd}
\text{(5)}\quad \mathcal{M}_{0,3}^{\sim} \arrow[rr] \arrow[dr] & & \mathrm{Aut}(\widehat{\mathcal{G}}^{+}) \arrow[dl] \\
& \mathrm{Aut}(\widehat{\mathcal{G}}_{0,3}^{+}=\widehat{\mathcal{G}}^{+}/\{\pm1\}) &
\end{tikzcd}}

This shows that (3) is injective, since
$\mathcal{M}_{0,3}^{\sim}\longrightarrow\mathrm{Aut}(\widehat{\mathcal{G}}_{0,3}^{+})$
is. \deleted{The map}

% source: bottom portion, lines 17--20, remark continuing on scan574.
\uline{Remark} The map
$\mathrm{Aut}(\widehat{\mathcal{G}}^{+})\longrightarrow\mathrm{Aut}(\mathcal{G}_{0,3}^{+\wedge})$
is not injective on the other hand, \deleted{\unsure{the}} the
elements of the kernel correspond to the
\unsure{homomorphisms}
$\widehat{\mathcal{G}}^{+}\longrightarrow\unsure{\mathrm{Centr}|\mathrm{Centre}}(\widehat{\mathcal{G}}^{+})=\{\pm1\}$,

\mspage{572}

% Printed archive footer outside the manuscript frame (not manuscript):
% 572
% https://grothendieck.umontpellier.fr
% Université de Montpellier, UMR CNRS 5149, Case courrier 051, Place Eugène Bataillon
% 34095 Montpellier cedex 5 - France

\scanpage{574}
% source: upper portion, lines 1--5, (6)--(7), including right marginal map.

or again to the homomorphisms \unsure{from}

\diagram{\begin{tikzcd}
\text{(6)}\quad \mathcal{G}_{\mathrm{ab}}^{+\wedge} \arrow[r] & \{\pm1\} & \text{or again}\quad (\mathcal{G}_{\mathrm{ab}}^{+})_{2} \arrow[r] & \{\pm1\}
\end{tikzcd}}

Now we see that
\[
\mathcal{G}_{\mathrm{ab}}^{+\wedge}
=\left\{\rho,\sigma\ \middle|\ \begin{gathered}
\rho^{6}=\sigma^{4}=[\rho,\sigma]=1\\
\rho^{3}=\sigma^{2}
\end{gathered}\right\}=\mathbb{Z}/12\mathbb{Z}
\]

\margin{\diagram{\begin{tikzcd}
\rho \arrow[r,mapsto] & 2\bmod12\\
\sigma \arrow[r,mapsto] & 3\bmod12
\end{tikzcd}}}

hence

\diagram{\begin{tikzcd}
& (\mathcal{G}_{\mathrm{ab}}^{+})_{2} \arrow[r,"\sim"] & \mathbb{Z}/2\mathbb{Z}\\
\text{(7)}\quad\left\{\vphantom{\begin{matrix}1\\1\\1\end{matrix}}\right. & \rho \arrow[r,mapsto] & 0\\
& \sigma \arrow[r,mapsto] & 1
\end{tikzcd}}

% source: middle portion, lines 6--9, composite diagram (8).
\deleted{hence to}
\inserted{This isomorphism corresponds to the character}
$\mathcal{G}^{+\wedge}\longrightarrow\{\pm1\}$ which is also obtained as the
composite

\diagram{\begin{tikzcd}[column sep=large]
\text{(8)}\quad \mathcal{G}^{+\wedge} \arrow[r] \arrow[rr,bend left=27,shorten >=50pt,"{\unsure{\mathrm{sg}|\mathrm{sgn}}}"] &
\underset{\substack{\text{symmetric group}\\\text{on three letters}}}{\mathcal{G}_{3}}
=\mathcal{G}^{+\wedge}/\underset{\substack{\text{inverse}\\\text{image}\\\text{of }\Pi_{0,3}^{\wedge}\\\text{in }\mathcal{G}^{+\wedge}}}{\widehat{\Pi}}
\arrow[r,"{\substack{\text{sign}\\\text{\unsure{of a} permutation}}}"] & \{\pm1\}
\end{tikzcd}}

% source: middle-lower portion, lines 10--13, (9)--(10).
Thus the only element \inserted{$\Theta$} \deleted{\unsure{different}} from $1$
of the kernel of \deleted{(3) is given by} \inserted{(4) is given by}

\diagram{\begin{tikzcd}
\text{(9)}\quad \deleted{\mathrm{Aut}\mathcal{G}}\ \inserted{\Theta:}\ x \arrow[r,mapsto] & x\cdot\mathrm{sg}(x)
\end{tikzcd}}

We shall note, \deleted{that \unsure{$\Theta$ and}} on setting
\[
(10)\qquad L_{0}^{\mathcal{G}}=\varepsilon_{0}^{\widehat{\mathbb{Z}}},\qquad
L_{0}^{\Pi}=L_{0}^{\mathcal{G}}\cap\Pi=l_{0}^{\widehat{\mathbb{Z}}},
\]
% source: lower portion, heavily rewritten lines 14--17, (11), ending unfinished.
that the subgroup $\Theta(L_{0}^{\mathcal{G}})$\inserted{$\subset\mathcal{G}^{+\wedge}$}
is \uline{not} $\mathcal{G}$-conjugate to $L_{0}$,
\deleted{\unsure{or}} \deleted{and $\Theta(L_{0}^{\Pi})$ $\mathcal{G}$-conjugate}
\inserted{for} \inserted{it is generated by}
\deleted{\unsure{$\Theta\varepsilon$}}\deleted{and a fortiori not.} :
\[
(11)\qquad \Theta(\varepsilon_{0})=-\varepsilon_{0}
\qquad\left(\begin{gathered}
\text{for }\varepsilon_{0}=\sigma\rho\\
\mathrm{sg}(\varepsilon_{0})=-1
\end{gathered}\right)
\]
and we check that $\varepsilon_{0}$ is not conjugate to
$(-\varepsilon_{0})^{\mu}=\deleted{\unsure{\varepsilon_{0}}}-\varepsilon_{0}^{\mu}$
for \unsure{any} $\mu\in\widehat{\mathbb{Z}}^{*}$ (as we see

\mspage{573}

% Printed archive footer outside the manuscript frame (not manuscript):
% 573
% https://grothendieck.umontpellier.fr
% Université de Montpellier, UMR CNRS 5149, Case courrier 051, Place Eugène Bataillon
% 34095 Montpellier cedex 5 - France

\scanpage{575}
\mspage{591}
% source: upper portion, pale continuation, lines 1--9 and (12).

at once, on \deleted{passing: $\mathcal{G}/\Pi_{\unsure{a}}$}\inserted{passing:}
$\mathcal{G}_{0,3}^{\wedge}/\Pi_{0,3}^{\wedge}{}_{\mathrm{ab}}$
and \unsure{noting} that $l_{0}=\varepsilon_{0}^{2}$ would be
conjugate to: $(-\varepsilon_{0}^{\mu})^{2}=l_{0}^{\mu}$,
\unsure{that we must have|which must have from it} $\mu=1$, now
$\varepsilon_{0}$ and $-\varepsilon_{0}$ are not conjugate
even in $\mathrm{Sl}(2,\widehat{\mathbb{Z}})$\unsure{'|!} ... ). Now we shall
be interested \unsure{precisely} in the automorphisms \unsure{of|$\varphi$}
$\mathcal{G}^{+\wedge}$ which transform $L_{0}^{\mathcal{G}}$ into a
conjugate subgroup. We \unsure{shall note|note} that on the other hand
\[
(12).\qquad \Theta(l_{0})=l_{0}\qquad\text{hence}\qquad
\Theta(L_{0}^{\Pi})=L_{0}^{\Pi}
\]

% source: middle portion, proposition and conditions a)--b), lines 10--17.
\uline{Proposition} The homomorphism (3)

\diagram{\begin{tikzcd}
\mathcal{M}_{0,3}^{\sim} \arrow[r] & \mathrm{Aut}(\widehat{\mathcal{G}}^{+})
\end{tikzcd}}

is injective, and its image consists of the
automorphisms\inserted{$u$} of $\widehat{\mathcal{G}}^{+}$ which satisfy the following
two conditions

a) $u$ stabilizes $\widehat{\Pi}=\mathcal{G}^{+\wedge}[2]
=\mathrm{Ker}(\mathcal{G}^{+\wedge}\longrightarrow\mathrm{Sl}(2,\mathbb{Z}/2\mathbb{Z}))$

b) $u$ transforms $L_{0}^{\mathcal{G}}$ into a subgroup of $\mathcal{G}^{+\wedge}$
$\mathcal{G}^{+\wedge}$-conjugate to $L_{0}$; i.e.
$\exists\mu\in\widehat{\mathbb{Z}}^{*}$ and $f_{0}\in\mathcal{G}^{+\wedge}$
such that we have
\[
(13)\qquad u(\varepsilon_{0})=\mathrm{int}(f_{0}^{-1})(\varepsilon_{0}^{\mu}).
\]

% source: lower portion, proof opening, lines 18--21.
\uline{Proof.} Injectivity has already been seen.
\deleted{\unsure{Let us prove|Let us recall}}\inserted{Let us show}
that conditions a)b) are necessary for
$u$ to \deleted{\unsure{belong}}\inserted{come}
from a $u_{0}\in\mathcal{M}_{0,3}^{\sim}$,
\deleted{For} a) this is trivial, since $u_{0}$ stabilizes
$\Pi_{0,3}^{\unsure{\wedge|\sim}}$, hence its

\mspage{574}

% Printed archive footer outside the manuscript frame (not manuscript):
% 574
% https://grothendieck.umontpellier.fr
% Université de Montpellier, UMR CNRS 5149, Case courrier 051, Place Eugène Bataillon
% 34095 Montpellier cedex 5 - France